\section*{\texorpdfstring{\S\CurrentExerciseGroup}{Section \CurrentExerciseGroup}}
\phantomsection
\addcontentsline{toc}{section}{\texorpdfstring{Exercises for \S\CurrentExerciseGroup}{Exercises for Section \CurrentExerciseGroup}}

\begin{enumerate}
\def\labelenumi{\arabic{enumi})}
\item
  Let E be a Hilbert space with a countable orthonormal basis \((\mathbf{e}_{n})\) (TVS, V, §2, No.3), \(E_{0}\) the dense linear subspace of E formed by the linear combinations of finitely many vectors \(e_{n}\) . Let X be the compact subspace of R formed by 0 and the points 1/n (n an integer \(\geqslant 1\) ); let \(\mu\) be the positive measure on X defined by the mass \(1/n^{2}\) at each point 1/n and the mass 0 at the point 0 (§1, No.~3, Example I). Consider the continuous mapping f of X into \(E_{0}\) defined by \(\mathbf{f}(0)=0\) , \(\mathbf{f}(1/n)=\mathbf{e}_{n}/n\) for \(n\geqslant1\) ; show that the integral \(\int f d\mu\) does not belong to \(E_{0}\) .
\item
  The spaces X and E and the mapping f being defined as in Exer. 1, show that the mapping \(\lambda \mapsto \lambda(\mathbf{f})\) of \(\mathcal{M}(\mathrm{X};\mathbf{C})\) into E is not continuous for the vague topology. (If \(g_{k}\) ( \(1 \leqslant k \leqslant n\) ) are \(n\) continuous scalar functions on X, show that there exists a discrete measure \(\lambda\) on X such that \(\int g_{k} d\lambda = 0\) for \(1 \leqslant k \leqslant n\) but \(\| \int \mathbf{f} d\lambda \|\) is arbitrarily large.)
\item
  \begin{enumerate}
  \def\labelenumii{\alph{enumii})}
  \tightlist
  \item
    Let E be a quasi-complete Hausdorff locally convex space. Show that when \(\mathcal{M}(\mathrm{X};\mathbf{C})\) is equipped with the quasi-strong topology (§1, Exer. 8) (resp. the topology of strictly compact convergence), the mapping \((\mathbf{f},\mu)\mapsto\int\mathbf{f}d\mu\) is \((\mathfrak{S},\mathfrak{T})\) -hypocontinuous, \(\mathfrak{T}\) being the set of vaguely bounded subsets of \(\mathcal{M}(\mathrm{X};\mathbf{C})\) , \(\mathfrak{S}\) the set of bounded (resp. compact) subsets of \(\mathcal{K}(\mathrm{X};\mathrm{E})\) that are contained in \(\mathcal{K}(\mathrm{X},\mathrm{K};\mathrm{E})\) for a suitable compact set K.
  \end{enumerate}
\end{enumerate}

\begin{enumerate}
\def\labelenumi{\alph{enumi})}
\setcounter{enumi}{1}
\tightlist
\item
  For every compact subset \(\mathbf{K}\) of \(\mathbf{X}\) , the mapping \((\mathbf{f},\mu)\mapsto \int \mathbf{f}d\mu\) is continuous on \(\mathcal{H}(\mathrm{X},\mathrm{K};\mathrm{E})\times \mathcal{M}(\mathrm{X};\mathbf{C})\) when \(\mathcal{M}(\mathrm{X};\mathbf{C})\) is equipped with the quasi-strong topology, and on \(\mathcal{H}(\mathrm{X},\mathrm{K};\mathrm{E})\times \mathcal{M}_{+}(\mathrm{X})\) when \(\mathcal{M}_{+}(\mathrm{X})\) is equipped with the vague topology.
\end{enumerate}

\begin{enumerate}
\def\labelenumi{\arabic{enumi})}
\setcounter{enumi}{3}
\tightlist
\item
  Prove the Cor. of Prop. 4 of No.~2 by making use of Prop. 8 of No.~4 (reduce to the case that E is quasi-complete, and make use also of §1, No.~9, Cor. 3 of Prop. 15).
\end{enumerate}

\setcounter{exercisegroup}{4}
\edef\CurrentExerciseGroup{\arabic{exercisegroup}}
